\chapter*{Introduction générale}
\addcontentsline{toc}{chapter}{Introduction générale}
\markboth{Introduction générale}{Introduction générale}

L'Internet des Objets (IoT) connaît une expansion exponentielle dans le secteur des télécommunications. Selon les estimations de la GSMA, le nombre de connexions IoT dépassera 25~milliards d'unités d'ici 2025, générant un volume de données sans précédent. Les opérateurs télécom comme \TT{} se trouvent au cœur de cette transformation, en déployant des infrastructures capables d'absorber ces flux tout en garantissant leur confidentialité, leur intégrité et leur disponibilité. Or, l'IoT présente des caractéristiques propres qui compliquent drastiquement la sécurisation : ressources calculatoires limitées, mises à jour rares, et protocoles applicatifs (MQTT, CoAP) conçus initialement sans mécanismes de sécurité obligatoires. Ces facteurs font des écosystèmes IoT des cibles de choix pour les attaquants, qui peuvent y déployer des \textit{botnets}, des attaques par déni de service, des interceptions \textit{Man-in-the-Middle} ou des injections de données malveillantes.

C'est dans ce cadre que s'inscrit notre projet de fin d'études, réalisé au sein de \TT{}, qui répond à un besoin concret et stratégique : définir, implémenter et valider une architecture de sécurité complète pour des flux MQTT en environnement IoT simulé, depuis l'authentification cryptographique des dispositifs jusqu'à la détection automatisée des comportements anormaux par Machine Learning.

\paragraph{Structure du rapport.}
Le présent document est organisé en cinq chapitres. Le \textbf{chapitre~1} introduit le contexte du projet, présente l'organisme d'accueil, formule la problématique et détaille la méthodologie Scrum adoptée ainsi que la planification temporelle. Le \textbf{chapitre~2} analyse les besoins fonctionnels et non fonctionnels et justifie les choix technologiques retenus à travers un état de l'art ciblé. Le \textbf{chapitre~3} décrit la réalisation et l'implémentation de l'architecture sécurisée (topologie réseau, PKI HashiCorp Vault, broker Mosquitto en TLS~1.3/mTLS). Le \textbf{chapitre~4} couvre la simulation de la flotte IoT et l'analyse réseau via Suricata et Wazuh. Le \textbf{chapitre~5} présente la détection d'anomalies par Machine Learning et l'évaluation quantitative des performances. Une conclusion générale dresse le bilan des objectifs atteints et propose des perspectives d'évolution. Deux annexes regroupent l'intégralité des scripts et configurations techniques.
